\section{The H/X physical allocation}

The source register order is $(N,T,W,A,D,H,X,Y,L)$.  Write
$p=(6,2,0,4,3,7,8,5,1)$ for the primary machine's logical-to-physical
map, and put $p'=(6,2,0,4,3,8,7,5,1)$.  Thus the new machine exchanges
the cached theorem register $H$ and scratch register $X$ on the tape.
It does not exchange their logical roles.

A physical field of value $v$ is $0\,1^{v+1}$.  The exchange acts on the
ordered list of field values, not by a fixed coordinate permutation of the
binary tape: unequal fields have unequal lengths.  This distinction matters
for the proof.  The source-address bits, header padding and nonempty spare
list are unchanged.  Let $F_p(s)$ and $F_{p'}(s)$ denote the two complete
physical frames for the same logical source configuration $s$.

Only two native transition rows change:
\[
 (21,1):0R60\longmapsto0R55,\qquad
 (45,1):0R55\longmapsto0R60.
\]
Controls $55$ and $60$ select physical field indices 7 and 8,
respectively (indices start at zero).  Every arithmetic-engine row remains identical to the primary
table.  The source graph still has 3,673 instruction sites.

\paragraph{Universal local refinement.}
The literal symbolic replay rederives four list-sweep lemmas and checks all
4,780 instruction outcomes, including both decrement branches.  Its affine
variables range over independent arbitrary natural field values and header
padding; the spare tail is an arbitrary nonempty finite list.  For each
nonhalting source step $s\longmapsto s'$, the new table executes a finite
positive-duration trace from $F_{p'}(s)$ to $F_{p'}(s')$, and no intermediate
control is halting.  For a halting source instruction, its trace reaches the
halt control.  The four generic list lemmas are proved from the literal
transition rows; the replay does not appeal to a register-arithmetic oracle.

The checked durations have nonnegative affine coefficients and constant at
least three.  The completed replay uses 1,068,062 symbolic operations, at
most 392 for an outcome.  These operation counts are costs of symbolic
verification, not physical machine runtimes.  The exact startup trace is
56,134 physical transitions.  Its endpoint is the whole source frame at
site 25, native control 84, head $-2$, logical vector $(2,0,\ldots,0)$ and
26 one-valued spares.  Since $H=X=0$ there, its binary tape is identical to
the primary machine's 58,154-step startup endpoint.  The saving is exactly
2,020 initialization transitions; no uniform later speedup is asserted.

\paragraph{Relative halting equivalence.}
The frame exchange and the positive-duration source contracts show that the
primary machine and PA194HX halt from blank tape under the same circumstances.
For one direction, concatenate the source instruction traces through the
first halting source instruction.  For reflection, a physical halt at time
$t$ cannot lie outside all initial source segments: if the source has not
halted after $t+1$ instructions, their positive durations cover more than
$t$ physical steps and none of those covered configurations halts.  This
contradicts the assumed physical halt.  The other table satisfies the same
source specification, so the argument applies in both directions.

\paragraph{Arithmetic proof boundary.}
This is a relative simulation result.  A weak-arithmetic implementation can
use primitive-recursive source iteration, accumulated segment lengths, and
bounded search through the first $t+1$ source steps to transport halting
witnesses.  If the local trace contracts and their composition are internalized
in iPRA, the same argument supplies an iPRA proof of the relative machine
halting equivalence.  This package supplies the explicit finite certificates
and the mathematical argument, not a newly compiled iPRA proof of that
internalization.  In particular it supplies no unconditional identification
with PA refutations.  It does not resolve the missing baseline calibration
\(\mathrm{iPRA}\vdash\operatorname{Con}(\mathrm{PA})\leftrightarrow
\neg\operatorname{Halt}(M)\).
